% Generated draft: requires mathematical and source review.
\documentclass[11pt]{article}
\usepackage[margin=25mm]{geometry}
\usepackage{fontspec}
\setmainfont{Latin Modern Roman}
\usepackage{amsmath,amssymb,mathtools,unicode-math}
\setmathfont{Latin Modern Math}
\usepackage{enumitem,needspace,xurl,hyperref}
\setlength{\parindent}{0pt}
\setlength{\parskip}{5pt}
\setlength{\emergencystretch}{3em}
\newcommand{\N}{\mathbb{N}}
\newcommand{\Z}{\mathbb{Z}}
\newcommand{\Q}{\mathbb{Q}}
\newcommand{\R}{\mathbb{R}}
\newcommand{\C}{\mathbb{C}}
\newcommand{\F}{\mathbb{F}}
\newcommand{\A}{\mathbb{A}}
\newcommand{\PP}{\mathbb{P}}
\newcommand{\E}{\mathbb{E}}
\newcommand{\eps}{\varepsilon}
\newcommand{\dd}{\,\mathrm{d}}
\begin{document}
\section*{UnsolvedMath ID: 30000461}
\textbf{Source problem code:} catalogue entry
\subsection{Definitions and mathematical statement}
\textbf{Domino-cell conjecture, as supplied.} In the type-\(B_n\) unequal-parameter Hecke setting, the partition into left cells should be given by domino insertion relative to an \(r\)-sized \(2\)-core, and hence should depend only on \(r\). The source does not state the parameter interval defining \(r\), what \emph{size} means, or which equivalence on domino tableaux describes a cell. These are material formulation gaps, especially at parameter walls.\par In a conventional setup, \(W(B_n)\) is the signed permutation group with simple generators \(t,s_1,\ldots,s_{n-1}\). A positive weight function has \(L(t)=b\) and \(L(s_i)=a\), with \(a,b>0\). Its Hecke algebra has basis \(T_w\) and the normalization \((T_s-v^{L(s)})(T_s+v^{-L(s)})=0\). The Kazhdan--Lusztig left preorder is generated by occurrences of basis elements in left multiplication in the Kazhdan--Lusztig basis, and its mutual-preorder equivalence classes are the left cells. The actual weight domain, normalization, and identification of cells in the source must be fixed, rather than assumed from the two parameter names alone.\par Standard domino insertion for signed permutations produces a pair of standard domino tableaux of the same shape, with a fixed \(2\)-core, each domino carrying a label. A \(2\)-core is a partition with no removable rim domino; the usual rank-\(d\) core is the staircase \(\delta_d=(d,d-1,\ldots,1)\), of \(d(d+1)/2\) boxes. Thus a core of rank \(r\) is not generally a core with \(r\) boxes, and a literal arbitrary box count need not occur for a \(2\)-core. The phrase \(r\)-sized is retained subject to recovering the intended convention. With rank \(d\) insertion, the tableaux have \(n\) dominos outside \(\delta_d\).\par A precise conjecture must give a rule assigning the appropriate core parameter to \(a,b\), specify the insertion and recording tableaux, and state whether left equivalence means equality of recording tableaux or an equivalence generated by specified allowed moves, such as cycle moves. Left versus right cells must use the correct insertion convention. In particular, integral or boundary values of \(b/a\) must not silently be treated by an unspecified rule for an open parameter interval. Dependence only on \(r\) means constancy of the cell partition within the relevant parameter class, not an undefined assertion that all unequal parameters have the same cells.\par The supplied triage checked 2026-08-21 labels the item partially solved. It reports proofs in equal-parameter, asymptotic, and intermediate regimes, including a case described as \(n-2\), but says its arbitrary-rank formulation is not fully settled. The exact hypotheses and meaning of \(n-2\) are not reproduced, and no blanket theorem is inferred from these names. This dated assessment requires independent verification.
\subsection{Short English statement}
The conjecture proposes a domino-tableau description of unequal-parameter type-\(B\) left cells, constant in the intended core parameter. The weight-to-core rule and tableau equivalence are missing. Core rank must be distinguished from its number of boxes, and parameter-wall cases need separate specification.
\subsection{Sources}
\begin{itemize}\item Local UnsolvedMath registry, permanent ID 30000461, code catalogue entry; \emph{Domino-Insertion Description of Hecke Left Cells}, Oberwolfach Reports Open Problems.\item \emph{Representations of Finite Groups (2007)}, Oberwolfach Reports, DOI 10.4171/owr/2006/15, pp. 895--978. \url{https://doi.org/10.4171/owr/2006/15}. Proposed year and year tag: 2006; no proposer named.\item Supplied literature label \emph{Howse--Pietraho, intermediate parameters}: \url{https://arxiv.org/abs/1902.09301}.\item Supplied literature label \emph{original Oberwolfach report}: \url{https://doi.org/10.4171/owr/2006/15}.\item Literature triage checked 2026-08-21, status \texttt{partially\_solved}; record updated 2026-09-16. The \texttt{corrected\_verified} field records notation restoration and explicit conjectural force, not a definition of the missing parameter rule. Verify dated claims independently.\end{itemize}
\end{document}
